%----------------------------------------------------------------------------------------
%	PACKAGES AND THEMES
%----------------------------------------------------------------------------------------
\documentclass[aspectratio=169,xcolor=dvipsnames]{beamer}
\usetheme{SimplePlus}
\usecolortheme{default}
\definecolor{grayedout}{RGB}{190,190,210}  % Choose whichever option you prefer

\usepackage{threeparttable}
\usepackage{hyperref}
\usepackage{subfig}
\usepackage{bbm}
\usepackage{graphicx} % Allows including images
\usepackage{booktabs} % Allows the use of \toprule, \midrule and \bottomrule in tables

\useoutertheme{beamerouterthemeprogressbar}
% Customize the colors (optional)
\setbeamercolor{progress bar}{fg=structure!50, bg=grayedout!30}
\setbeamercolor{progress bar in head/foot}{fg=structure!50, bg=grayedout!30}
\setbeamercolor{progress bar in section page}{fg=structure!50, bg=grayedout!30}
%----------------------------------------------------------------------------------------
%	TITLE PAGE
%----------------------------------------------------------------------------------------

\title[short title]{ICFES Referrals} % The short title appears at the bottom of every slide, the full title is only on the title page
\subtitle{10-minute internal presentation for data analysis}

\author[Pin-Yen] {Reha Tuncer - University of Luxembourg}

\date{27 March 2025} % Date, can be changed to a custom date


%----------------------------------------------------------------------------------------
%	PRESENTATION SLIDES
%----------------------------------------------------------------------------------------

\begin{document}

\begin{frame}
    % Print the title page as the first slide
    \titlepage
\end{frame}

\begin{frame}{Selection into the experiment}
    \begin{columns}[T]  % T option aligns the top of the columns
        \begin{column}{0.38\textwidth}  % Left column for your comment
            \vspace{2em}  % Add some space to align better with the table
            \begin{itemize}
                \item Better students self select
                \item Low and High SES self select
            \end{itemize}
        \end{column}
        \begin{column}{0.62\textwidth}  % Right column for the table
            \begin{table}
                \centering
                \begin{threeparttable}
                \begin{tabular*}{0.9\columnwidth}{@{\extracolsep{\fill}}lcc c@{}}
                \toprule
                & \multicolumn{1}{c}{\textbf{Admin Data}} & \multicolumn{1}{c}{\textbf{Sample}} & \multicolumn{1}{c}{\textbf{p}} \\
                \midrule
                Reading score & 62.651 & 65.183 & 0.000 \\
                Math score & 63.973 & 67.477 & 0.000 \\
                GPA & 3.958 & 4.012 & 0.000 \\
                Low-SES & 0.343 & 0.410 & 0.000 \\
                Med-SES & 0.505 & 0.499 & 0.763 \\
                High-SES & 0.153 & 0.091 & 0.000 \\
                \midrule
                Observations & 4,417 & 734 & 5,151 \\
                \bottomrule
                \end{tabular*}
                \label{tab:selection}
                \end{threeparttable}
            \end{table}
        \end{column}
    \end{columns}
\end{frame}

\begin{frame}{Balance between treatments}
    \begin{columns}[T]  % T option aligns the top of the columns
        \begin{column}{0.38\textwidth}  % Left column for placeholder text
            \vspace{2em}  % Add some space to align better with the table
            \begin{itemize}
                \item Successful randomization
            \end{itemize}
        \end{column}
        \begin{column}{0.62\textwidth}  % Right column for the table
            \begin{table}
                \centering
                \begin{threeparttable}
                \begin{tabular*}{0.9\columnwidth}{@{\extracolsep{\fill}}lcc c@{}}
                \toprule
                & \multicolumn{1}{c}{\textbf{Baseline}} & \multicolumn{1}{c}{\textbf{Bonus}} & \multicolumn{1}{c}{\textbf{p}} \\
                \midrule
                Reading score & 64.712 & 65.693 & 0.134 \\
                Math score & 67.366 & 67.597 & 0.780 \\
                GPA & 4.003 & 4.021 & 0.445 \\
                \# connections & 173.40 & 176.88 & 0.574 \\
                Tie strength & 3.939 & 3.719 & 0.443 \\
                Low-SES & 0.419 & 0.401 & 0.615 \\
                Med-SES & 0.492 & 0.506 & 0.714 \\
                High-SES & 0.089 & 0.094 & 0.824 \\
                \midrule
                Observations & 382 & 352 & 734 \\
                \bottomrule
                \end{tabular*}
                \label{tab:balance}
                \end{threeparttable}
            \end{table}
        \end{column}
    \end{columns}
\end{frame}

\begin{frame}{Network size and tie strength by semester}
    \begin{columns}[T]
        \begin{column}{0.38\textwidth}
            \vspace{2em}
            \begin{itemize}
                \item Classes together increases over time
                \item Connections peak around 7 semesters and decline as students graduate
            \end{itemize}
        \end{column}
        \begin{column}{0.62\textwidth}
            \centering
            \includegraphics[width=0.99\columnwidth]{/Users/reha.tuncer/Documents/GitHub/icfes-referrals/figures/connection_class.png}
            \label{fig:connections_class}
        \end{column}
    \end{columns}
\end{frame}

\begin{frame}{Referrer network composition by SES}
    \begin{columns}[T]
        \begin{column}{0.38\textwidth}
            \vspace{2em}
            \begin{itemize}
                \item Monotonic increase in availability across SES
                \item All differences are statistically significant ($p < 0.001$)
            \end{itemize}
        \end{column}
        \begin{column}{0.62\textwidth}
            \centering
            \includegraphics[width=0.95\columnwidth]{/Users/reha.tuncer/Documents/GitHub/icfes-referrals/figures/ses_distribution.png}
            \label{fig:baseline_ses}
        \end{column}
    \end{columns}
\end{frame}

\begin{frame}{Referrer network performance by SES}
    \begin{columns}[T]
        \begin{column}{0.38\textwidth}
            \vspace{2em}
            \begin{itemize}
                \item Subject specific network-level standardization
                \item Performance increases with SES across math and reading
                \item All differences are statistically significant ($p < 0.005$)
            \end{itemize}
        \end{column}
        \begin{column}{0.62\textwidth}
            \centering
            \includegraphics[width=0.95\columnwidth]{/Users/reha.tuncer/Documents/GitHub/icfes-referrals/figures/ses_peer_performance.png}
            \label{fig:performance}
        \end{column}
    \end{columns}
\end{frame}

\begin{frame}{Referrals for Reading}
    \begin{columns}[T]
        \begin{column}{0.38\textwidth}
            \vspace{2em}
            \begin{itemize}
                \item Referrals have higher reading scores and much higher tie strength on aggregate (all $p < 0.001$)
                \item No treatment effect on reading score or on on tie strength (all $p > 0.08$)
            \end{itemize}
        \end{column}
        \begin{column}{0.62\textwidth}
            \centering
            \includegraphics[width=0.95\columnwidth]{/Users/reha.tuncer/Documents/GitHub/icfes-referrals/figures/reading_combined.png}
            \label{fig:reading_combined}
        \end{column}
    \end{columns}
\end{frame}

\begin{frame}{Referrals for Math}
    \begin{columns}[T]
        \begin{column}{0.38\textwidth}
            \vspace{2em}
            \begin{itemize}
                \item Referrals have higher math scores and much higher tie strength on aggregate (all $p < 0.001$)
                \item No treatment effect on math score or on on tie strength (all $p > 0.1$)
            \end{itemize}
        \end{column}
        \begin{column}{0.62\textwidth}
            \centering
            \includegraphics[width=0.95\columnwidth]{/Users/reha.tuncer/Documents/GitHub/icfes-referrals/figures/math_combined.png}
            \label{fig:math_combined}
        \end{column}
    \end{columns}
\end{frame}

\begin{frame}{Baseline referrals}
    \centering
    \begin{columns}
        \begin{column}{0.45\textwidth}
            \centering
            \includegraphics[width=\textwidth]{/Users/reha.tuncer/Documents/GitHub/icfes-referrals/figures/ses_referral_distribution.png}
            \label{fig:baseline_rates}
        \end{column}
        \hfill
        \begin{column}{0.45\textwidth}
            \centering
            \includegraphics[width=\textwidth]{/Users/reha.tuncer/Documents/GitHub/icfes-referrals/figures/baseline_performance.png}
            \label{fig:baseline_performance}
        \end{column}
    \end{columns}

    \begin{itemize}
        \item Large monotonic differences in referral shares as SES increases (all $p < 0.1$)
        \item Small differences in math and reading performance across SES (all $p > 0.08$)
    \end{itemize}
\end{frame}

\begin{frame}{Treatment effects on referrals}
    \centering
    \begin{columns}
        \begin{column}{0.45\textwidth}
            \centering
            \includegraphics[width=\textwidth]{/Users/reha.tuncer/Documents/GitHub/icfes-referrals/figures/ses_treatment_effects.png}
            \label{fig:bonus_rates}
        \end{column}
        \hfill
        \begin{column}{0.45\textwidth}
            \centering
            \includegraphics[width=\textwidth]{/Users/reha.tuncer/Documents/GitHub/icfes-referrals/figures/treatment_effect.png}
            \label{fig:bonus_performance}
        \end{column}
    \end{columns}

    \begin{itemize}
        \item No meaningful changes in referral shares with bonus
        \item No meaningful changes in reading and math performance with bonus
    \end{itemize}
\end{frame}

\begin{frame}{Treatment effect on referral tie strength}
    \centering
    \begin{columns}
        \begin{column}{0.45\textwidth}
            \centering
            \includegraphics[width=\textwidth]{/Users/reha.tuncer/Documents/GitHub/icfes-referrals/figures/baseline_ties.png}
            \label{fig:baseline_tie}
        \end{column}
        \hfill
        \begin{column}{0.45\textwidth}
            \centering
            \includegraphics[width=\textwidth]{/Users/reha.tuncer/Documents/GitHub/icfes-referrals/figures/ties_treatment_effects.png}
            \label{fig:bonus_tie}
        \end{column}
    \end{columns}

    \begin{itemize}
        \item Small differences in referral tie strength at baseline
        \item No meaningful changes in referral tie strength with bonus
    \end{itemize}
\end{frame}

\begin{frame}[plain]
    \vspace{3em}
    \begin{center}
        \Huge\textbf{Regression Analysis}
    \end{center}
\end{frame}

\begin{frame}{Empirical specification}
    \begin{itemize}
        \item Model the referral choice using a mixed multinomial logit model for each referrer's choice set (average 175 potential nominees) [McFadden and Train, 2000]
        
        \item The probability that student $i$ refers student $j \in N_i$ is given by
        $$\frac{e^{\beta_i x_{ij}}}{\sum_{j=1}^{N_i} e^{\beta_i x_{ij}}}$$
        
        where $x_{ij}$ is the vector of characteristics of $j$ (low-SES, test score, tie strength) and $\beta_i$ is a vector of individual-specific parameters comprising
        of a constant and a random component
    \end{itemize}
\end{frame}

\begin{frame}{Is there a Low-SES bias in Reading referrals?}
    \begin{columns}[T]
        \begin{column}{0.38\textwidth}
            \vspace{2em}
            \begin{itemize}
                \item Reading score and tie strength are strong predictors of referrals
                \item No interaction between reading score and tie strength
                \item No evidence for a Low-SES bias in referrals when controlling for referral score and tie strength
            \end{itemize}
        \end{column}
        \begin{column}{0.62\textwidth}
            \begin{table}
                \centering
                \footnotesize
                \begin{threeparttable}
                \begin{tabular*}{0.9\columnwidth}{@{\extracolsep{\fill}}lccc}
                \toprule
                 & \multicolumn{1}{c}{(1)} & \multicolumn{1}{c}{(2)} & \multicolumn{1}{c}{(3)} \\
                \midrule
                Low-SES & 0.199** & 0.041 & 0.042 \\
                 & (0.083) & (0.100) & (0.100) \\
                Reading Score  &  & 0.561*** & 0.509*** \\
                 &  & (0.044) & (0.048) \\
                Tie &  & 0.951*** & 0.941*** \\
                 &  & (0.031) & (0.032) \\
                Score x Tie &  &  & 0.029 \\
                 &  &  & (0.018) \\
                \midrule
                Observations & 128,847 & 128,847 & 128,847 \\
                Individuals & 673 & 673 & 673 \\
                \bottomrule
                \end{tabular*}
                \label{tab:reading_regression}
                \end{threeparttable}
            \end{table}
        \end{column}
    \end{columns}
\end{frame}

\begin{frame}{Is there a Low-SES bias in Math referrals?}
    \begin{columns}[T]
        \begin{column}{0.38\textwidth}
            \vspace{2em}
            \begin{itemize}
                \item Math score and tie strength are strong predictors of referrals
                \item Significant but small interaction between math score and tie strength
                \item No evidence for a Low-SES bias in referrals when controlling for referral score and tie strength
            \end{itemize}
        \end{column}
        \begin{column}{0.62\textwidth}
            \begin{table}
                \centering
                \footnotesize
                \begin{threeparttable}
                \begin{tabular*}{0.9\columnwidth}{@{\extracolsep{\fill}}lccc}
                \toprule
                 & \multicolumn{1}{c}{(1)} & \multicolumn{1}{c}{(2)} & \multicolumn{1}{c}{(3)} \\
                \midrule
                Low-SES & 0.220*** & 0.049 & 0.050 \\
                 & (0.083) & (0.097) & (0.098) \\
                Math Score  &  & 0.653*** & 0.538*** \\
                 &  & (0.040) & (0.041) \\
                Tie &  & 0.887*** & 0.854*** \\
                 &  & (0.029) & (0.030) \\
                Score x Tie &  &  & 0.088*** \\
                 &  &  & (0.019) \\
                \midrule
                Observations & 128,150 & 128,150 & 128,150 \\
                Individuals & 669 & 669 & 669 \\
                \bottomrule
                \end{tabular*}
                \label{tab:math_regression}
                \end{threeparttable}
            \end{table}
        \end{column}
    \end{columns}
\end{frame}

\begin{frame}{Do grades predict Reading referrals?}
    \begin{columns}[T]
        \begin{column}{0.38\textwidth}
            \vspace{2em}
            \begin{itemize}
                \item GPA correlates weakly with reading score ($\rho = 0.18$)
                \item GPA is a better predictor of referrals when controlling for reading score and tie strength
            \end{itemize}
        \end{column}
        \begin{column}{0.62\textwidth}
            \begin{table}
                \centering
                \footnotesize
                \begin{threeparttable}
                \begin{tabular*}{0.9\columnwidth}{@{\extracolsep{\fill}}lccc}
                \toprule
                 & \multicolumn{1}{c}{(1)} & \multicolumn{1}{c}{(2)} & \multicolumn{1}{c}{(4)} \\
                \midrule
                GPA & 0.759*** & 0.640*** & 0.578*** \\
                 & (0.044) & (0.059) & (0.070) \\
                Reading Score &  & 0.407*** & 0.344*** \\
                 &  & (0.048) & (0.068) \\
                Tie Strength &  & 0.936*** & 0.911*** \\
                 &  & (0.031) & (0.032) \\
                Score × Tie &  &  & 0.010 \\
                 &  &  & (0.022) \\
                Score × GPA &  &  & 0.017 \\
                 &  &  & (0.056) \\
                GPA × Tie &  &  & 0.028 \\
                 &  &  & (0.026) \\
                Score × Tie × GPA &  &  & 0.030* \\
                 &  &  & (0.018) \\
                \bottomrule
                \end{tabular*}
                \label{tab:reading_gpa_regression}
                \end{threeparttable}
            \end{table}
        \end{column}
    \end{columns}
\end{frame}

\begin{frame}{Do grades predict Math referrals?}
    \begin{columns}[T]
        \begin{column}{0.38\textwidth}
            \vspace{2em}
            \begin{itemize}
                \item Significant but small interactions between tie strength and both performance variables
                \item GPA correlates weakly with math score ($\rho = 0.14$)
                \item GPA is a better predictor of referrals when controlling for math score and tie strength
            \end{itemize}
        \end{column}
        \begin{column}{0.62\textwidth}
            \begin{table}
                \centering
                \footnotesize
                \begin{threeparttable}
                \begin{tabular*}{0.9\columnwidth}{@{\extracolsep{\fill}}lccc}
                \toprule
                 & \multicolumn{1}{c}{(1)} & \multicolumn{1}{c}{(2)} & \multicolumn{1}{c}{(4)} \\
                \midrule
                GPA & 0.765*** & 0.621*** & 0.500*** \\
                 & (0.044) & (0.059) & (0.068) \\
                Math Score &  & 0.531*** & 0.400*** \\
                 &  & (0.044) & (0.065) \\
                Tie Strength &  & 0.871*** & 0.808*** \\
                 &  & (0.029) & (0.030) \\
                Score × Tie &  &  & 0.068*** \\
                 &  &  & (0.023) \\
                Score × GPA &  &  & 0.026 \\
                 &  &  & (0.061) \\
                GPA × Tie &  &  & 0.069*** \\
                 &  &  & (0.026) \\
                Score × Tie × GPA &  &  & 0.034 \\
                 &  &  & (0.022) \\
                \bottomrule
                \end{tabular*}
                \label{tab:math_gpa_regression}
                \end{threeparttable}
            \end{table}
        \end{column}
    \end{columns}
\end{frame}

\begin{frame}{Is there in-group homophily in Reading referrals?}
    \begin{columns}[T]
        \begin{column}{0.38\textwidth}
            \vspace{2em}
            \begin{itemize}
                \item Restrict sample by SES of referrer
                \item Define homophily as same SES referrer-nominee dyad
                \item Only Low-SES exhibit homophily for other Low-SES
            \end{itemize}
        \end{column}
        \begin{column}{0.62\textwidth}
            \begin{table}
                \centering
                \footnotesize
                \begin{threeparttable}
                \begin{tabular*}{0.9\columnwidth}{@{\extracolsep{\fill}}lccc}
                \toprule
                & Low-SES & Med-SES & High-SES \\
                 & \multicolumn{1}{c}{(1)} & \multicolumn{1}{c}{(2)} & \multicolumn{1}{c}{(3)} \\
                \midrule
                Homophily & 0.300** & 0.203 & -0.354 \\
                 & (0.152) & (0.125) & (0.344) \\
                Reading Score  & 0.433*** & 0.353*** & 0.674*** \\
                 & (0.074) & (0.067) & (0.170) \\
                Tie & 0.860*** & 0.991*** & 1.061*** \\
                 & (0.043) & (0.048) & (0.134) \\
                GPA & 0.648*** & 0.667*** & 0.477** \\
                 & (0.094) & (0.081) & (0.228) \\
                \midrule
                Observations & 54,611 & 64,596 & 9,640 \\
                Individuals & 275 & 340 & 58 \\
                \bottomrule
                \end{tabular*}
                \label{tab:homophily_reading}
                \end{threeparttable}
            \end{table}
        \end{column}
    \end{columns}
\end{frame}

\begin{frame}{Is there in-group homophily in Math referrals?}
    \begin{columns}[T]
        \begin{column}{0.38\textwidth}
            \vspace{2em}
            \begin{itemize}
                \item Restrict sample by SES of referrer
                \item Define homophily as same SES referrer-nominee dyad
                \item Only Low-SES exhibit homophily for other Low-SES
            \end{itemize}
        \end{column}
        \begin{column}{0.62\textwidth}
            \begin{table}
                \centering
                \footnotesize
                \begin{threeparttable}
                \begin{tabular*}{0.9\columnwidth}{@{\extracolsep{\fill}}lccc}
                \toprule
                & Low-SES & Med-SES & High-SES \\
                 & \multicolumn{1}{c}{(1)} & \multicolumn{1}{c}{(2)} & \multicolumn{1}{c}{(3)} \\
                \midrule
                Homophily & 0.282* & 0.161 & -0.030 \\
                 & (0.146) & (0.125) & (0.358) \\
                Math Score  & 0.500*** & 0.516*** & 0.810*** \\
                 & (0.069) & (0.064) & (0.128) \\
                Tie & 0.823*** & 0.892*** & 1.046*** \\
                 & (0.040) & (0.044) & (0.141) \\
                GPA & 0.580*** & 0.675*** & 0.519** \\
                 & (0.093) & (0.083) & (0.219) \\
                \midrule
                Observations & 55,531 & 62,492 & 10,127 \\
                Individuals & 283 & 327 & 59 \\
                \bottomrule
                \end{tabular*}
                \label{tab:homophily_math}
                \end{threeparttable}
            \end{table}
        \end{column}
    \end{columns}
\end{frame}

\begin{frame}{Empirical Specification}
    
    \vspace{0.5em}
    \textbf{OLS Model:}
    \begin{align*}
        \text{Score}_{ij} = \beta_0 &+ \beta_1 \cdot \text{X}_i + \beta_2 \cdot \text{Mean Score}_j + \beta_3 \cdot \text{SD Score}_j + \varepsilon_{ij}
    \end{align*}
    
    \vspace{0.5em}
    \begin{itemize}
        \item $\text{Score}_{ij}$: Test score of the nominee from referrer $i$ in subject $j$
        \item $\text{X}_{ij}$: SES of student $i$ or the test score of student $i$ in subject $j$
        \item $\text{Mean Score}_j$ and $\text{SD Score}_j$: Mean and SD of $i$'s network performance in subject $j$
    \end{itemize}
    
    \vspace{0.5em}
    \textbf{Features:}
    \begin{itemize}
        \item Sample restricted to nominated students only
        \item SES as categorical variable
        \item Controls for network-level performance differences
        \item Standard errors clustered at individual level
    \end{itemize}
\end{frame}

\begin{frame}{Do Low-SES refer better?}
    \begin{columns}[T]
        \begin{column}{0.42\textwidth}
            \vspace{2em}
            \begin{itemize}
                \item No evidence that Low-SES refer better in reading or math
                \item Better referrals depend on average network performance
            \end{itemize}
        \end{column}
        \begin{column}{0.58\textwidth}
            \begin{table}
                \centering
                \footnotesize
                \begin{tabular}{lcc}
                \toprule
                 & \multicolumn{1}{c}{Reading} & \multicolumn{1}{c}{Math} \\
                 & \multicolumn{1}{c}{(1)} & \multicolumn{1}{c}{(2)} \\
                \midrule
                Med-SES & -1.213** & -0.816 \\
                 & (0.607) & (0.731) \\
                High-SES & 0.438 & 0.937 \\
                 & (1.085) & (1.287) \\
                Mean Score & 1.250*** & 1.399*** \\
                 & (0.096) & (0.074) \\
                SD Score & 0.233 & 0.163 \\
                 & (0.409) & (0.318) \\
                Constant & -13.062 & -21.571*** \\
                 & (8.006) & (6.373) \\
                \midrule
                Observations & 673 & 669 \\
                \bottomrule
                \end{tabular}
            \end{table}
        \end{column}
    \end{columns}
\end{frame}

\begin{frame}{Is there performance homophily in referrals?}
    \begin{columns}[T]
        \begin{column}{0.42\textwidth}
            \vspace{2em}
            \begin{itemize}
                \item Significant but small effect of own test score on referral score
                \item Very large difference between the average network test score versus referrer performance (about 6 times)
                \item Performance homophily is nearly not as important as network composition
            \end{itemize}
        \end{column}
        \begin{column}{0.58\textwidth}
            \begin{table}
                \centering
                \footnotesize
                \begin{tabular}{lcc}
                \toprule
                 & \multicolumn{1}{c}{Reading} & \multicolumn{1}{c}{Math} \\
                 & \multicolumn{1}{c}{(1)} & \multicolumn{1}{c}{(2)} \\
                \midrule
                Own Score & 0.177*** & 0.176*** \\
                 & (0.035) & (0.036) \\
                Mean Score & 1.020*** & 1.186*** \\ 
                 & (0.101) & (0.082) \\
                SD Score & 0.056 & 0.155 \\
                 & (0.398) & (0.310) \\
                Constant & -9.126 & -19.835*** \\
                 & (7.835) & (6.240) \\
                \midrule
                Observations & 673 & 669 \\
                \bottomrule
                \end{tabular}
            \end{table}
        \end{column}
    \end{columns}
\end{frame}

\begin{frame}{Prelude}
    \begin{columns}[T]
        \begin{column}{0.38\textwidth}
            \vspace{2em}
            \begin{itemize}
                \item No evidence for a Low-SES bias in referrals
                \item Referrers pick nominees based on performance and tie strength
                \item Referral performance depends heavily on referrer network composition
                \item Which SES group has the better network?
            \end{itemize}
        \end{column}
        \begin{column}{0.62\textwidth}
            \centering
            \includegraphics[width=0.9\columnwidth]{/Users/reha.tuncer/Documents/GitHub/icfes-referrals/figures/tie_combined.png}
        \end{column}
    \end{columns}
\end{frame}

\end{document}
